\section{The inherited exact interface and new proof obligations}
\label{sec:interface}

\subsection{Input, rings, and immutable entries}

The input is a bipartite Boolean tensor network with at most $s$ wires,
polynomially many vertices, maximum degree at most $\lceil s^{9/10}\rceil$,
and description length $\mathcal L$. Its contraction is
\[
 Z=\sum_{x\in\{0,1\}^{E}}\prod_v T_v(x|_{\delta(v)}).
\]
Entries lie in $(\mathbb Z/2^{P_0}\mathbb Z)[V]/(V^\ell)$ with
$1\le P_0,\ell\le s$. A uniform entry routine takes
$(s+\mathcal L)^{O(1)}2^{O(s^{9/10})}$ time and polynomial space, with
fixed constants. The low-degree theorem changes the degree premise alone.

The source uses an extension ring
\[
 R_p=(\mathbb Z/2^p\mathbb Z)[U,V]/(f(U),V^\ell)
\]
with polynomial rank and precision. Each tensor is tagged with its creation
precision. Later evaluations use the same canonical integer representatives
of its creation-time value. Canonical lifting is not treated as a ring
homomorphism. These conventions, the deformation, and exact Hermite recovery
are imported \cite{source,asymmetric,joint,flags,trees}.

\subsection{Exact mode sum and entropy}

At an active call with $M$ wires, fixed alphabet $r$, retained taper $u$,
packet exponent $a$, and valuation $i<u$, put $D=2^{a-i}$. The compatible
inherited lemmas give child wire count $(2/r+o(1))M$, precision
$p_t\le(6+o(1))MD/u$, and at most $(2/r^2+o(1))M/D$ mode families of
each index on each directed pair. Their certified weights are
$a+1-i+4j-2\pc(j)$.

Enumerate an exact weight allocation $(b_{gh})_{g\ne h}$ with sum below
$p_t$. Let $\mathcal I_{gh}$ be its finite list of mode tuples. Skip an
allocation with an empty list and set
\[
 h_{gh}=\log_2|\mathcal I_{gh}|,\qquad
 \K(M,u)=\kappa M/\sqrt u,\qquad\kappa=\pi\sqrt2/\ln2.
\]
Its contribution is exactly
\begin{equation}\label{eq:exact-sum}
 \sum_{\xi_{gh}\in\mathcal I_{gh}}
       \prod_{g=1}^r Z_g((\xi_e)_{e\ni g}),
\end{equation}
where each $Z_g$ entry is a genuine recursive contraction. There are
polynomially many allocations since $r$ is fixed and all budgets are
polynomial in $s$. A knapsack recurrence computes exact list cardinalities,
rank/unrank, and polynomial-delay iteration; no counting oracle is introduced.

For every fixed nonnegative directed vector $z$, the common-tilt estimate is
\begin{equation}\label{eq:entropy}
 \sum_{g\ne h}z_{gh}h_{gh}
 \le \left(\frac{\norm z}{r}+o(1)\right)\K(M,u).
\end{equation}
To recall its leading constant, the mode weight generating product has
natural logarithm $\pi^2/(24\theta)+O(\log^2(2/\theta))$.
The shared budget then gives the following natural-logarithm bound,
before division by $\ln2$:
\[
 \theta p_t+\frac{\pi^2}{24\theta}
        \frac{(2+o(1))M}{r^2D}\sum z_{gh}^2.
\]
Optimizing the first two terms gives \eqref{eq:entropy}; the additive buffer
makes the error uniform on the fixed finite support family used here.

For unordered pairs $e=\{g,h\}$ write $g_e=h_{gh}+h_{hg}$.
Any unordered support $v$ therefore satisfies
\begin{equation}\label{eq:undirected-entropy}
 v\cdot g\le\left(\frac{\sqrt2\norm v}{r}+o(1)\right)\K.
\end{equation}
The factor $\sqrt2$ is essential. A fixed finite-net argument also gives
$\norm h\le(1+o(1))\K/r$ when needed below.

\subsection{Packing gates and the paid recurrence}

With $L=2^a$, $Q=r^a$, side degrees $R,C$, sufficient gates are
\begin{equation}\label{eq:gates}
 QR/M=o(1),\qquad Q^2C(\log s)^2/M=o(1),\qquad LQ^2/M=o(1).
\end{equation}
Initially both degrees are at most $s^{9/10+o(1)}$; afterwards the retained
side remains under that bound and the newly created side has degree at most
$2^{a_{d-1}+1}$. The low-degree proof tracks all previous packet degrees.

Each group has $2(r-1)$ incident directed pairs. Thus forming its full table
requires at most $2^{(\alpha_r+o(1))\K}$ recursive child evaluations, where
\[
 \alpha_r=\frac{\sqrt{2(r-1)}}r.
\]
If a certified local contraction plan costs $2^{(\beta_r+o(1))\K}$,
the correctly charged recurrence is
\begin{equation}\label{eq:recurrence}
 T_d\le\poly(s)\left[
    2^{(\alpha_{r_d}+o(1))\K_d}T_{d+1}
    +2^{(\beta_{r_d}+o(1))\K_d}\right].
\end{equation}
A smaller local operation coefficient alone is insufficient; the theorem
must control the maximum of all branches of this recurrence.

\subsection{Rectangular backend retained from the corpus}

We use the already reviewed field bounds $\alpha>3/10$ and $\omega<19/8$,
supported by the stronger primary results in \cite{matrix}. Fix
$\varepsilon=10^{-5}$. For sorted logarithmic matrix dimensions
$x\ge y\ge z\ge0$, the corpus's rectangular construction gives
\begin{equation}\label{eq:rectangle}
 F(x,y,z)=\max\left\{
 (1+\varepsilon/2)(x+y),\quad
 x+(47/56+\varepsilon)y+(15/28)z\right\}.
\end{equation}
The first term pays for input/output arrays. It cannot be omitted.
Tensoring a fixed $N\times N^{3/10}\times N$ algorithm of exponent
$2+\varepsilon$ with a fixed square algorithm of exponent $19/8$ gives
the second term, with square blocking of the largest dimension. Rational
coordinate lifting, nonexceptional odd primes, and CRT extend the algorithms
to $R_p$ without dividing by a nonunit. Fixed bilinear recursions process one
child product at a time and use a polynomial factor times the largest
input/output array in space \cite{rectangular}.

For completeness, the rectangle interpolation takes
$t_0=(y-z)/(1-3/10)$ and $t_1=(z-(3/10)y)/(1-3/10)$ when
$z\ge(3/10)y$. The arithmetic exponent, including blocking, is bounded by
$x-y+(2+\varepsilon)t_0+(19/8)t_1$.
When $z<(3/10)y$, padding gives $x+y+\varepsilon y$.
Both are at most \eqref{eq:rectangle}. The three matrix dimension
permutations are justified by the symmetries of the bilinear multiplication
tensor, with each permuted algorithm executed afresh.

The present contribution changes the adaptive plan certificate. It does not
claim the rectangular matrix bound, Cauchy basis, or Hermite recovery as new.
